\chapter{Avviare e fermare il programma}
\label{ch:launcher}
\index{Avvio}\index{Launcher}

\lettrine[lines=3]{P}{er} usare piTantum non serve digitare comandi
complicati: c'\`e uno script che avvia tutto e uno che ferma tutto.
piTantum \`e composto da due pezzi che partono insieme --- il
\emph{motore} (backend) e l'\emph{interfaccia} nel browser
(frontend) --- ma dal tuo punto di vista \`e un'unica applicazione.

\section{Avviare}

\begin{description}
\item[Windows] Fai doppio click su \texttt{webui\textbackslash
start.bat} (oppure eseguilo da terminale). Si aprono due finestre
--- lasciale aperte: sono il motore e l'interfaccia in esecuzione.
\item[Linux / macOS] Da terminale, nella cartella del progetto:
\texttt{webui/start.sh}. Parte in background; aggiungi
\texttt{-{}-foreground} se preferisci tenerlo in primo piano e
fermarlo con \texttt{Ctrl+C}.
\end{description}

Quando \`e pronto, apri il browser su \texttt{http://localhost:5173}:
\`e l'indirizzo dell'interfaccia.

\begin{avvertenzabox}
\textbf{Il primo avvio richiede qualche minuto}: la prima volta lo
script prepara da solo l'ambiente Python e le dipendenze del
frontend. Dalla seconda volta in poi l'avvio \`e immediato. Se una
delle due porte (\texttt{8000} per il motore, \texttt{5173} per
l'interfaccia) risulta gi\`a occupata, lo script te lo segnala:
probabilmente \`e rimasta aperta una sessione precedente ---
fermala prima (vedi sotto).
\end{avvertenzabox}

\section{Fermare}

\begin{description}
\item[Windows] Chiudi le due finestre aperte all'avvio.
\item[Linux / macOS] \texttt{webui/stop.sh}: ferma in modo pulito
sia il motore sia l'interfaccia.
\end{description}

\section{Cosa fanno gli script, in breve}

Entrambi gli avvii, prima di partire, fanno alcuni controlli:
verificano che Python e Node siano disponibili, creano l'ambiente
virtuale in \texttt{backend/.venv} se manca, installano le
dipendenze del frontend al primo giro, e avvisano se le porte
\texttt{8000}/\texttt{5173} sono gi\`a in uso. Su Linux/macOS
suggeriscono anche il comando giusto (\texttt{apt}, \texttt{dnf},
\texttt{pacman}, \texttt{brew}) se manca qualcosa.

\begin{biblioparvabox}
\textbf{Dettagli tecnici.} Su Windows \texttt{start.bat} usa path
relativi a \texttt{\%\textasciitilde dp0}, quindi funziona ovunque
lo si copi. Su Linux/macOS \texttt{start.sh} scrive i log in
\texttt{webui/logs/} e i PID in \texttt{webui/logs/pids};
\texttt{stop.sh} li rilegge, invia \texttt{SIGTERM} con fallback
\texttt{SIGKILL} dopo 3\,s, termina anche i processi figli
(\texttt{pkill -P}) e, come ultima rete, libera le porte
\texttt{8000}/\texttt{5173} via \texttt{lsof}.
\end{biblioparvabox}
